problem:  
Let $(M^n,g)$ be a compact Riemannian manifold with **harmonic curvature**, i.e. $\nabla^*{\mathrm{Riem}}=0$, but not locally symmetric and not Einstein.  By the second Bianchi identity, this is equivalent to the Ricci tensor being Codazzi.  In all known examples (notably those classified by Derdziński in dimension four and certain warped products in higher dimension), such metrics exhibit a local warped‑product or conformally Einstein structure.  

**Question:**  
In dimensions $n\ge 5$, does there exist a compact, non‑conformally‑Einstein, non‑product Riemannian manifold $(M^n,g)$ with harmonic curvature and irreducible universal cover?  
Equivalently, can harmonic curvature in dimensions $\ge5$ occur without any underlying warped‑product or conformal Einstein structure—i.e. are Derdziński‑type examples the only possible local models, or is there genuinely new geometric behavior at higher dimension?  

Why is it a "good" problem:  
This formulation removes redundant assumptions and targets the core unknown: the global and local geometry of higher‑dimensional harmonic‑curvature manifolds.  The question probes whether the analytic condition $\nabla^*\mathrm{Riem}=0$ enforces an inherent decomposition or conformal rigidity beyond the known four‑dimensional theory.  

Resolving it would determine whether harmonic‑curvature metrics in higher dimensions form a finite, rigid family (Einstein, product, or conformally Einstein) or whether they constitute a broader, flexible geometry yet to be classified.  The problem bridges Riemannian analysis, global structure theorems, and conformal geometry, and could reveal new geometric phenomena paralleling the leap from four to higher dimensional conformal structures—making it a concise yet profound research‑level problem.